% Pista pro-26s-q3 · Probabilitat
% Procedència: PAU Catalunya 2026 · Convocatòria de setembre · Sèrie 2 · Exercici 3
\documentclass[11pt,a4paper]{article}
\usepackage{pista}
\pistainfo{Catalunya 2026}{Convocatòria de setembre}{Sèrie 2}

\begin{document}

\section*{\textcolor{blauPista}{Exercici 3}}

\noindent Es vol construir un parc eòlic al Solsonès format per 10 aerogeneradors. L'empresa que els construirà està dissenyant les pales segons l'esquema següent, on $x$ i $f(x)$ estan mesurats en metres (el dibuix no està a escala):

\noindent\textit{(A la figura: $g(x) = 8x$ i $f(x) = 2x^{3}$; la part ombrejada és tota la regió compresa entre les dues gràfiques.)}

\noindent S'ha calculat que cadascun dels 10 aerogeneradors tindrà una potència de $100\,\text{MW}$. Degut a diversos factors, cadascun té una probabilitat del $10\,\%$ d'estar aturat, independentment de la resta d'aerogeneradors.

\subsection*{3.1. \normalfont Quina és la probabilitat que s'arribi als $1\,000\,\text{MW}$ de potència total al parc? \hspace*{\fill}\textnormal{[0,75~punts]}}

\pista{Anomena $X$ el nombre d'aerogeneradors en funcionament. Com que són 10 i cadascun funciona o no independentment dels altres, $X \sim B(10, p)$, on $p$ és la probabilitat que funcioni cada aerogenerador.
\begin{itemize}
\item Quants aerogeneradors han de funcionar per arribar als $1\,000\,\text{MW}$?
\item Calcula la probabilitat d'aquest únic cas.
\end{itemize}

\medskip
\textbf{Atenció:} el $10\,\%$ és la probabilitat d'estar \emph{aturat}; la de funcionar és la complementària.}

\subsection*{3.2. \normalfont Quina és la probabilitat que la potència total del parc superi els $150\,\text{MW}$? \hspace*{\fill}\textnormal{[0,75~punts]}}

\pista{Cada aerogenerador aporta $100\,\text{MW}$ o res.
\begin{itemize}
\item Tradueix «superar els $150\,\text{MW}$» a una condició sobre $X$: quin és el nombre mínim d'aerogeneradors en funcionament?
\item Sumar tots els casos favorables seria llarguíssim: passa a l'esdeveniment contrari, que només té dos casos.
\item Quan en funciona exactament un, pot ser qualsevol dels 10 (el coeficient binomial).
\end{itemize}

\medskip
\textbf{Atenció:} les probabilitats del contrari són molt petites: escriu-les en notació científica i no arrodoneixis el resultat final.}

\subsection*{3.3. \normalfont Calculeu l'àrea de les pales d'un d'aquests aerogeneradors (part ombrejada al gràfic). \hspace*{\fill}\textnormal{[1~punt]}}

\pista{Calcula l'àrea en dos passos:
\begin{itemize}
\item resol $f(x) = g(x)$: els punts de tall són els límits d'integració;
\item a cada interval entre dos punts de tall consecutius, calcula el valor absolut de la integral de la diferència,
\[
\left| \int_{x_{i}}^{x_{i+1}} \bigl(f(x) - g(x)\bigr)\,dx \right|,
\]
i suma els resultats (recorda les unitats, $\text{m}^{2}$).
\end{itemize}

\medskip
\textbf{Atenció:} si integres d'un extrem a l'altre d'una sola vegada, les dues regions es compensen.

\medskip
\textbf{Drecera:} les dues funcions són senars (simètriques respecte de l'origen): calcula una sola regió i multiplica per 2.}

\end{document}
